%============================================================
% CH2 · PATTERN 3 — ENTRY AND EXIT TIMES: THEORY
%============================================================
\subsection{Pattern Theory}
\label{ssec:ch2-p3-theory}

% the running tree T of this pattern: per-vertex styles can be overridden in a scope
\tikzset{
    g2p3-v1/.style=gv, g2p3-v2/.style=gv, g2p3-v3/.style=gv, g2p3-v4/.style=gv, g2p3-v5/.style=gv,
    g2p3-v6/.style=gv, g2p3-v7/.style=gv, g2p3-v8/.style=gv, g2p3-v9/.style=gv,
    g2p3-e/.style=gtree,
    g2p3-T/.pic={
        \node[g2p3-v1] (t1) at (3,0) {1};
        \node[g2p3-v2] (t2) at (1.2,-1.2) {2};
        \node[g2p3-v3] (t3) at (4.8,-1.2) {3};
        \node[g2p3-v4] (t4) at (0.4,-2.4) {4};
        \node[g2p3-v5] (t5) at (2.0,-2.4) {5};
        \node[g2p3-v6] (t6) at (4.0,-2.4) {6};
        \node[g2p3-v7] (t7) at (5.6,-2.4) {7};
        \node[g2p3-v8] (t8) at (3.4,-3.6) {8};
        \node[g2p3-v9] (t9) at (4.6,-3.6) {9};
        \foreach \a/\b in {1/2, 1/3, 2/4, 2/5, 3/6, 3/7, 6/8, 6/9} \draw[g2p3-e] (t\a) -- (t\b);
    }
}

\begin{patternbox}{Pattern 3 at a Glance}
Run \textbf{one} DFS and write down, for every vertex, \emph{when} it was entered and \emph{when} it was left. Afterwards the tree is no longer needed: ``$v$ is below $u$'' is two comparisons, ``the subtree of $u$'' is one contiguous block of an array, ``the descendants of $u$ at depth $d$'' is a binary search, and the order in which vertices \emph{finish} is an order in which they can be peeled off without disconnecting the rest.

\begin{center}
\begin{tikzpicture}[x=0.9cm, y=0.9cm]
    \tikzset{gv/.append style={minimum size=6.5mm, font=\scriptsize\bfseries}}
    \begin{scope}[g2p3-v3/.style=gvcur, g2p3-v6/.style=gvseen, g2p3-v7/.style=gvseen,
                  g2p3-v8/.style=gvseen, g2p3-v9/.style=gvseen]
        \pic {g2p3-T};
    \end{scope}
    \foreach \v/\a/\b/\p in {1/0/17/above, 2/1/6/left, 3/7/16/right, 4/2/3/below, 5/4/5/below,
                             6/8/13/left, 7/14/15/below, 8/9/10/below, 9/11/12/below}
        \node[gdist, text=darkgray, \p=0.4mm of t\v] {\tiny[\a,\b]};
    % the pre-order array
    \begin{scope}[xshift=7.3cm, yshift=-1.0cm]
        \foreach \v/\i in {1/0, 2/1, 4/2, 5/3, 3/4, 6/5, 8/6, 9/7, 7/8}{
            \ifnum\i<4 \node[gcell, minimum size=6.5mm, font=\scriptsize\bfseries] (c\i) at (\i*0.75,0) {\v};
            \else \node[gcell, fill=hitfill, minimum size=6.5mm, font=\scriptsize\bfseries] (c\i) at (\i*0.75,0) {\v}; \fi
            \node[font=\tiny, text=darkgray] at (\i*0.75,-0.55) {\i};
        }
        \node[gcell, fill=currentfill, draw=accentorange, line width=1pt, minimum size=6.5mm, font=\scriptsize\bfseries] at (4*0.75,0) {3};
        \draw[accentorange, line width=1.2pt, decorate, decoration={brace, amplitude=4pt}] (2.75,0.45) -- (6.25,0.45)
            node[midway, above=4pt, font=\scriptsize, text=accentorange] {pos$[3] = 4$, sz$[3] = 5$};
        \node[font=\scriptsize, anchor=west] at (-0.4,-1.2) {pre-order \texttt{order[]} (index $=$ pos)};
        \node[font=\scriptsize, anchor=west, text width=6.4cm] at (-0.4,-2.0) {$3$ is above $9$: $7 \le 11$ and $12 \le 16$.\quad $2$ is not above $6$: $\tout{6} = 13 > 6$.};
    \end{scope}
\end{tikzpicture}
\end{center}
\begin{recurrencebox}
\centering\small
enter $u$: $\tin{u} = \texttt{timer++}$, $\mathrm{pos}[u] = \texttt{cnt++}$ \qquad leave $u$: $\tout{u} = \texttt{timer++}$, $\mathrm{sz}[u] = \texttt{cnt} - \mathrm{pos}[u]$, append $u$ to \texttt{post}\\[3pt]
$u$ above $v \iff \tin{u} \le \tin{v} \wedge \tout{v} \le \tout{u}$ \qquad subtree$(u) = \texttt{order}[\mathrm{pos}[u] \,..\, \mathrm{pos}[u]+\mathrm{sz}[u]-1]$\\[3pt]
one DFS \VE, then \bigO{1} per ancestor query and \bigO{\log n} per ``at depth $d$'' query
\end{recurrencebox}
\end{patternbox}

%------------------------------------------------------------
\subsubsection*{From the lessons to the pattern}

Lesson~\ref{ssec:ch2-timestamps} introduced the two numberings and proved the parenthesis theorem; Pattern~2 used DFS to \emph{compute} something for every subtree while the search was running. This pattern turns the point of view around: the DFS runs \textbf{once, first}, and records only times. All questions are answered \emph{afterwards}, from the arrays, often many thousands of them. The single new ingredient is the observation that time is an \emph{order}, and orders can be searched:
\begin{itemize}[leftmargin=*]
    \item the interval $[\tin{u}, \tout{u}]$ contains exactly the entry times of $u$'s descendants (parenthesis theorem), so ``descendant'' becomes ``number in an interval'';
    \item in pre-order numbering the descendants are not only inside an interval, they \emph{fill} it: the subtree is a block of consecutive positions;
    \item the finishing order (post-order) lists every vertex after all its descendants --- the order for peeling, and, on DAGs, the reverse of a topological order.
\end{itemize}

%------------------------------------------------------------
\subsubsection*{When to recognise it (and when not)}

\begin{itemize}[leftmargin=*]
    \item many queries ``is $u$ an ancestor of $v$?'', ``does the path from the root to $v$ pass through $u$?'', ``is $v$ in the subtree of $u$?'' --- the \bigO{1} test;
    \item ``the $k$-th vertex visited when a DFS starts at $u$'', ``list the subtree of $u$ in DFS order'' --- the pre-order block (Problem~\ref{ssec:ch2-military});
    \item ``how many descendants of $u$ have depth $d$ / colour $c$ / label $\ell$'' --- per group a sorted list of entry times and a binary search, or an offline counter difference (Problem~\ref{ssec:ch2-count-descendants});
    \item ``remove $k$ vertices (cells) so that the rest stays connected'', ``in which order can the pieces be dismantled'' --- the finishing order (Problem~\ref{ssec:ch2-maze});
    \item ``the vertices are numbered so that every subtree is consecutive'' --- the input \emph{is} a pre-order: use the block directly (E3.5).
\end{itemize}
\textbf{Not this pattern:} if a subtree value can be computed bottom-up during the DFS itself (sizes, sums, counts), Pattern~2 is simpler. If the queries \emph{change} values (``add 5 to vertex $v$, then ask the sum of a subtree''), the block is still the key, but the range structure on top of it (Fenwick tree, segment tree) belongs to Chapter~15. Lowest common ancestors and distances between two arbitrary vertices need binary lifting or an Euler tour with range minimum (Chapter~14). Ancestor tests in a general directed graph speak about the DFS tree, \emph{not} about reachability (Lesson~\ref{ssec:ch2-timestamps}).

%------------------------------------------------------------
\subsubsection*{The framework}

\begin{center}
\small
\renewcommand{\arraystretch}{1.3}
\begin{tabularx}{\linewidth}{@{}>{\bfseries}l X@{}}
\toprule
Record & on entry: $\tin{u}$, $\mathrm{pos}[u]$, $\mathrm{dep}[u]$, append $u$ to \texttt{order}; on exit: $\tout{u}$, $\mathrm{sz}[u]$, append $u$ to \texttt{post} \\
Ancestor & $\tin{u} \le \tin{v}$ and $\tout{v} \le \tout{u}$ (or with the block: $\mathrm{pos}[u] \le \mathrm{pos}[v] < \mathrm{pos}[u] + \mathrm{sz}[u]$) \\
Subtree list & $\texttt{order}[\mathrm{pos}[u]], \dots, \texttt{order}[\mathrm{pos}[u] + \mathrm{sz}[u] - 1]$; the first element is $u$ itself \\
Group count & \texttt{byDepth[d]} $=$ entry times of depth-$d$ vertices (sorted for free); count those in $[\tin{u}, \tout{u}]$ \\
Peel order & \texttt{post}: every vertex after all of its descendants \\
Cost & one DFS \VE{} (iterative), $\bigO{n}$ memory per array; queries \bigO{1} or \bigO{\log n} \\
\bottomrule
\end{tabularx}
\end{center}

\begin{invariantbox}[Why a subtree is a block, and why peeling in finishing order is safe]
\textbf{Invariant (pre-order).} While \texttt{dfs(u)} is open, every vertex that receives a position is a descendant of $u$, and every descendant of $u$ receives its position while \texttt{dfs(u)} is open (white-path theorem: at $\tin{u}$ all of $u$'s subtree is white). Positions are handed out by one counter, so the positions given during \texttt{dfs(u)} are the consecutive numbers $\mathrm{pos}[u], \mathrm{pos}[u]+1, \dots, \texttt{cnt}-1$ at the moment $u$ is left --- exactly $\mathrm{sz}[u]$ of them. That is the block.

\textbf{Invariant (post-order).} At any moment, the set $F$ of finished vertices is closed under taking descendants: a vertex finishes only after all its children did. Hence for every \emph{unfinished} vertex $v \ne$ root, its parent is unfinished too (a finished parent would have finished $v$ before). So the unfinished vertices, together with their tree edges, form a connected subtree containing the root. Removing the first $k$ vertices of \texttt{post} therefore never disconnects the rest --- in any connected graph, since the DFS tree spans it.

\textbf{Depth lists are sorted.} Vertices are appended to \texttt{byDepth[d]} at their entry time, which only grows; so each list is increasing and binary search applies with no sort. \hfill$\blacksquare$
\end{invariantbox}

\begin{diagrambox}
\begin{tikzpicture}[x=0.62cm, y=0.5cm]
    \draw[->, darkgray] (-0.3,0) -- (18.4,0) node[right, font=\scriptsize] {time};
    \foreach \t in {0,...,17} {\draw[darkgray] (\t,0.12) -- (\t,-0.12); \node[font=\tiny, below] at (\t,-0.12) {\t};}
    \foreach \v/\a/\b/\l/\c in {1/0/17/5/primaryblue, 2/1/6/4/successgreen, 3/7/16/4/successgreen,
                                4/2/3/3/accentorange, 5/4/5/3/accentorange, 6/8/13/3/accentorange, 7/14/15/3/accentorange,
                                8/9/10/2/softpurple, 9/11/12/2/softpurple}{
        \draw[\c, line width=2.2pt] (\a,\l*0.95) -- (\b,\l*0.95);
        \draw[\c, line width=1pt] (\a,\l*0.95-0.25) -- (\a,\l*0.95+0.25);
        \draw[\c, line width=1pt] (\b,\l*0.95-0.25) -- (\b,\l*0.95+0.25);
        \node[font=\scriptsize\bfseries, fill=white, inner sep=1pt, text=\c] at ({(\a+\b)/2},\l*0.95+0.45) {\v};
    }
    \fill[accentorange, opacity=0.12] (7,-0.9) rectangle (16,5.5);
    \node[font=\scriptsize, text=accentorange, anchor=north] at (11.5,-0.95) {window of 3: every interval inside it is a descendant of 3};
\end{tikzpicture}
\figcap{The intervals $[\tin{u}, \tout{u}]$ of $T$, one row per depth. The ancestor test asks ``is $v$'s interval inside $u$'s window?''; because intervals never overlap partially, it is enough to look at the entry time $\tin{v}$ alone: $\tin{u} \le \tin{v} \le \tout{u}$.}
\end{diagrambox}

%------------------------------------------------------------
\subsubsection*{Combining the block with groups: per-depth lists}

``Descendants of $u$ with depth $d$'' mixes a \emph{subtree} condition with a \emph{group} condition. Keep one sorted list per group (here: per depth) of entry times. The subtree condition becomes an interval of times, and an interval of a sorted list is counted by two binary searches:
\[
    \#\{v \text{ below } u,\ \mathrm{dep}[v] = d\} \;=\; \texttt{upper\_bound}(L_d, \tout{u}) - \texttt{lower\_bound}(L_d, \tin{u}).
\]

\begin{diagrambox}
\begin{tikzpicture}[x=1cm, y=0.8cm]
    \foreach \d/\lst in {0/{0/1}, 1/{1/2, 7/3}, 2/{2/4, 4/5, 8/6, 14/7}, 3/{9/8, 11/9}}{
        \node[font=\scriptsize\bfseries, anchor=east] at (-0.2,-\d) {$L_{\d}$};
        \foreach \t/\v [count=\i from 0] in \lst {
            \node[gcell, minimum width=9mm, minimum height=6mm, font=\scriptsize] (L\d-\i) at (\i*1.0+0.3,-\d) {\t};
            \node[font=\tiny, text=darkgray, anchor=west] at (\i*1.0+0.3+0.32,-\d-0.28) {v\v};
        }
    }
    \node[gcell, minimum width=9mm, minimum height=6mm, font=\scriptsize\bfseries, fill=answerfill, draw=accentorange] at (0.3,-3) {9};
    \node[gcell, minimum width=9mm, minimum height=6mm, font=\scriptsize\bfseries, fill=answerfill, draw=accentorange] at (1.3,-3) {11};
    \node[gcell, minimum width=9mm, minimum height=6mm, font=\scriptsize\bfseries, fill=answerfill, draw=accentorange] at (2.3,-2) {8};
    \node[gcell, minimum width=9mm, minimum height=6mm, font=\scriptsize\bfseries, fill=answerfill, draw=accentorange] at (3.3,-2) {14};
    \node[note, text width=6.8cm, align=left, anchor=west] at (4.6,-1.0) {query $(u = 3, d = 3)$: window $[7, 16]$; in $L_3 = (9, 11)$ the first entry $\ge 7$ has index 0, the first $> 16$ has index 2: answer $2$ (vertices 8, 9).\\[3pt] query $(u = 3, d = 2)$: in $L_2 = (2, 4, 8, 14)$ the window covers indices $2..3$: answer $2$ (vertices 6, 7).\\[3pt] query $(u = 4, d = 1)$: $d < \mathrm{dep}[4]$, answer $0$.};
\end{tikzpicture}
\figcap{Per-depth lists of entry times of $T$ (small labels: the vertex). Each list is filled in time order, hence sorted. Values from the template program below.}
\end{diagrambox}

The same trick works for any partition of the vertices into groups (colour, label, letter): one list per group, $n$ entries in total. When the queries can be collected first, an offline variant needs no binary search at all: keep a counter per group, read it when entering $u$ and again when leaving $u$; the difference counts the group inside the subtree (Approach~3 of Problem~\ref{ssec:ch2-count-descendants}).

%------------------------------------------------------------
\subsubsection*{The finishing order: peeling and topological order}

The post-order is the third array a DFS gives for free. Two facts make it useful:
\begin{enumerate}[leftmargin=*]
    \item \textbf{Peeling.} The first vertices to finish are leaves of the DFS tree; in general every prefix of \texttt{post} is a union of complete subtrees, so removing a prefix leaves a connected graph (invariant above). This is the whole solution of Problem~\ref{ssec:ch2-maze}. The mirror statement for the \emph{pre}-order (and for a BFS order) is: every prefix is connected, because each vertex is discovered from an earlier one.
    \item \textbf{Reverse topological order.} In a DAG every edge $u \to v$ satisfies $\tout{v} < \tout{u}$: $v$ is white (tree edge, finishes inside $u$), or black (already finished), and it can never be gray (that would be a back edge, a cycle). So listing the vertices by \emph{decreasing} exit time puts every edge forwards: a topological order. Chapter~5 builds on this; here it is one \texttt{reverse}.
\end{enumerate}

\begin{diagrambox}
\begin{tikzpicture}[x=0.9cm, y=0.9cm]
    \tikzset{gv/.append style={minimum size=6.5mm, font=\scriptsize\bfseries}}
    \begin{scope}[g2p3-v4/.style=gvdead, g2p3-v5/.style=gvdead, g2p3-v2/.style=gvdead, g2p3-v8/.style=gvdead]
        \pic {g2p3-T};
        \foreach \a/\b in {1/2, 2/4, 2/5, 6/8} \draw[draw=white, line width=3pt] (t\a) -- (t\b);
        \foreach \a/\b in {1/2, 2/4, 2/5, 6/8} \draw[gfaded, dashed] (t\a) -- (t\b);
        \node[font=\scriptsize\bfseries] at (3,0.9) {remove the first 4 of \texttt{post}};
        \node[font=\scriptsize\ttfamily] at (3,-4.5) {post = \textcolor{midgray}{4 5 2 8} 9 6 7 3 1};
    \end{scope}
    \begin{scope}[xshift=7.6cm, yshift=-0.6cm]
        \node[font=\scriptsize\bfseries] at (2.0,1.5) {a DAG and its finishing order};
        \node[gvdone] (a1) at (0,0) {1}; \node[gvdone] (a2) at (1.4,0.7) {2}; \node[gvdone] (a3) at (1.4,-0.7) {3};
        \node[gvdone] (a4) at (2.8,0) {4}; \node[gvdone] (a5) at (4.2,0) {5}; \node[gvdone] (a6) at (0,-1.6) {6};
        \foreach \a/\b in {1/2, 1/3, 2/4, 3/4, 4/5, 6/3} \draw[gde] (a\a) -- (a\b);
        \foreach \v/\f/\p in {1/4/left, 2/2/above, 3/3/below, 4/1/above, 5/0/above, 6/5/left}
            \node[gdist, text=deepblue, \p=0.5mm of a\v] {\tiny\#\f};
        \node[font=\scriptsize\ttfamily, anchor=west] at (-0.3,-2.6) {finish: 5 4 2 3 1 6};
        \node[font=\scriptsize\ttfamily, anchor=west] at (-0.3,-3.15) {reverse: 6 1 3 2 4 5};
        \node[note, anchor=west, text width=4.6cm, align=left] at (-0.3,-3.95) {every edge goes left to right in the reversed list: a topological order (Chapter~5)};
    \end{scope}
\end{tikzpicture}
\figcap{Left: the four earliest finishers of $T$ are removed; what remains ($1, 3, 6, 7, 9$) is still a tree. Right: \#$i$ $=$ position in the finishing order; reversing it orders the DAG. Both orders printed by the template program.}
\end{diagrambox}

%------------------------------------------------------------
\subsubsection*{Computing the times iteratively}

With $n = 2\cdot10^5$ and a tree that may be a path, a recursive DFS needs $2\cdot10^5$ frames --- far beyond the about $65{,}000$ that the default 2\,MB stack survives (Lesson~\ref{ssec:ch2-call-stack}). The iterative form keeps, for every vertex on the stack, the index \texttt{it[u]} of the next neighbour to try. ``Enter'' happens when a vertex is pushed, ``leave'' when its iterator runs out and it is popped --- exactly the two moments where the recursive version writes \texttt{tin} and \texttt{tout}:

\begin{diagrambox}
\begin{tikzpicture}[node distance=4mm]
    \node[fcproc=3.0cm] (top) {$u =$ top of stack};
    \node[fcdec=3.0cm, below=of top] (d) {\texttt{it[u]} $<$ \texttt{deg(u)}?};
    \node[fcproc=3.6cm, right=14mm of d] (enter) {$v = \texttt{adj[u][it[u]++]}$; if $v \ne$ parent: \textbf{enter} $v$ (tin, pos, dep), push};
    \node[fcproc=3.6cm, below=of d] (leave) {\textbf{leave} $u$: tout, $\mathrm{sz}[u] = \texttt{cnt} - \mathrm{pos}[u]$, \texttt{post}, pop};
    \draw[fcarrow] (top) -- (d);
    \draw[fcarrow] (d) -- node[lbl, above]{yes} (enter);
    \draw[fcarrow] (d) -- node[lbl, right]{no} (leave);
    \draw[fcarrow] (enter.north) |- (top.east);
    \draw[fcarrow] (leave.west) -- ++(-0.9,0) |- (top.west);
\end{tikzpicture}
\figcap{One step of the iterative timestamp DFS. Each vertex is entered once and left once; each adjacency entry is read once: \VE{} and no recursion.}
\end{diagrambox}

When the input gives parents with $p_i < i$ (Problems~\ref{ssec:ch2-military} and~\ref{ssec:ch2-count-descendants}), sizes and depths need no DFS at all: depths by a loop $i = 2..n$, sizes by a loop $i = n..2$. The pre-order still needs the children order, but even it can be computed without a stack by handing out blocks top-down (Approach~3 of Problem~\ref{ssec:ch2-military}).

%------------------------------------------------------------
\subsubsection*{Variations map}

\begin{diagrambox}
\begin{tikzpicture}[
    var/.style={rectangle, rounded corners=2pt, draw=patternborder, fill=patternbg,
                font=\scriptsize, align=center, text width=3.25cm, minimum height=1.3cm, inner sep=3pt}]
    \node[var, fill=answerfill, draw=accentorange, text width=5.4cm, font=\scriptsize\bfseries] (root) at (0,0)
        {one DFS records tin / tout / pos / sz / post; answer from the arrays};
    \node[var] (v1) at (-5.4,-2) {\textbf{$k$-th in DFS order}\\\texttt{order[pos[u]+k-1]}\\(Military, 2.3.1)};
    \node[var] (v2) at (-1.8,-2) {\textbf{count at depth $d$}\\sorted tin per depth\\(Count Descendants, 2.3.2)};
    \node[var] (v3) at (1.8,-2)  {\textbf{peel keeping connectivity}\\first $k$ of post-order\\(Maze, 2.3.3)};
    \node[var] (v4) at (5.4,-2)  {\textbf{ancestor queries}\\two comparisons\\(E3.1)};
    \node[var] (v5) at (-5.4,-3.6) {\textbf{Euler walk}\\print on entry and on return\\(E3.2)};
    \node[var] (v6) at (-1.8,-3.6) {\textbf{edge types after the fact}\\tree / back / forward / cross\\(E3.3)};
    \node[var] (v7) at (1.8,-3.6)  {\textbf{which branch holds $x$?}\\child $c$ with $c$ above $x$\\(E3.4)};
    \node[var] (v8) at (5.4,-3.6)  {\textbf{input already in pre-order}\\edges cross the block twice\\(E3.5)};
    \node[draw=patternborder, dashed, rounded corners=4pt, fit=(v1)(v4)(v5)(v8), inner sep=5pt] (grp) {};
    \draw[flow] (root.south) -- (grp.north);
    \node[note, below=2mm of grp] {range sums / updates on the block: Chapter 15 \quad flattening as a technique: Chapter 13 \quad LCA: Chapter 14};
\end{tikzpicture}
\end{diagrambox}

\textbf{Twists and how to handle them}
\begin{itemize}[leftmargin=*]
    \item \textbf{``Children in increasing order''.} The pre-order depends on the order of the adjacency lists. Sort them, or build them so that they come out sorted (with $p_i < i$, appending $i$ to \texttt{ch[p\_i]} for $i = 2..n$ already sorts them).
    \item \textbf{The root is not vertex 1.} Run the DFS from the given root; for a question about ``the side of $x$ as seen from $y$'', simply root at $y$ (E3.4).
    \item \textbf{Strict or non-strict ancestor?} The test with $\le$ counts $u = v$ as an ancestor. For ``proper'' add $u \ne v$.
    \item \textbf{Both conventions at once.} One clock ($0..2n-1$) is best for comparisons; the pre-order counter ($0..n-1$) is best for indexing arrays. Recording both costs nothing.
    \item \textbf{A forest or a disconnected graph.} Keep the timer running across the outer loop; intervals of different trees are disjoint, so the tests stay correct.
    \item \textbf{The depth may exceed the list count.} Guard \texttt{d < byDepth.size()} (and $d \ge \mathrm{dep}[u]$) before indexing.
\end{itemize}

%------------------------------------------------------------
\subsubsection*{Code templates}

\begin{lstlisting}[style=cpp]
#include <bits/stdc++.h>
using namespace std;

// (1) Rooted tree (undirected adjacency lists, 1-indexed), ITERATIVE: safe for n = 2e5 on a path.
//     One clock:   tin[u] < tout[u] in 0..2n-1           -> ancestor test
//     Pre-order:   pos[u] in 0..n-1, order[pos[u]] = u   -> subtree = order[pos[u] .. pos[u]+sz[u]-1]
//     Post-order:  post = vertices in the order they finish (leaves of the DFS tree come early)
struct Timestamps {
    vector<int> tin, tout, pos, order, sz, dep, par, post;
    vector<vector<int>> byDepth;                  // byDepth[d] = tin of all depth-d vertices, sorted

    Timestamps(const vector<vector<int>>& adj, int root) {
        int n = adj.size() - 1, timer = 0, cnt = 0;
        tin.assign(n + 1, -1); tout = pos = sz = dep = par = tin;
        order.assign(n, 0);
        vector<size_t> it(n + 1, 0);
        vector<int> st;
        auto enter = [&](int v, int p) {
            par[v] = p; dep[v] = (p ? dep[p] + 1 : 0);
            tin[v] = timer++; pos[v] = cnt; order[cnt++] = v;
            if ((int)byDepth.size() <= dep[v]) byDepth.push_back(vector<int>());
            byDepth[dep[v]].push_back(tin[v]);    // pushed in increasing tin: already sorted
            st.push_back(v);
        };
        enter(root, 0);
        while (!st.empty()) {
            int u = st.back();
            if (it[u] < adj[u].size()) {
                int v = adj[u][it[u]++];
                if (v != par[u]) enter(v, u);
            } else {                                // u finishes
                tout[u] = timer++;
                sz[u] = cnt - pos[u];               // vertices entered since u
                post.push_back(u);
                st.pop_back();
            }
        }
    }
    bool isAnc(int u, int v) const { return tin[u] <= tin[v] && tout[v] <= tout[u]; }
    int kth(int u, int k) const { return k <= sz[u] ? order[pos[u] + k - 1] : -1; }  // k-th in u's subtree
    int countAtDepth(int u, int d) const {          // descendants of u (u included) at depth d
        if (d >= (int)byDepth.size() || d < dep[u]) return 0;
        const vector<int>& L = byDepth[d];
        return upper_bound(L.begin(), L.end(), tout[u]) - lower_bound(L.begin(), L.end(), tin[u]);
    }
};

// (2) Finishing order of a DIRECTED graph (all start vertices), iterative.
//     On a DAG the reverse of this order is a topological order (Chapter 5).
vector<int> finishOrder(const vector<vector<int>>& adj) {
    int n = adj.size() - 1;
    vector<int> color(n + 1, 0), res, st;
    vector<size_t> it(n + 1, 0);
    for (int s = 1; s <= n; s++) {
        if (color[s]) continue;
        color[s] = 1; st.push_back(s);
        while (!st.empty()) {
            int u = st.back();
            if (it[u] < adj[u].size()) {
                int v = adj[u][it[u]++];
                if (!color[v]) { color[v] = 1; st.push_back(v); }
            } else { color[u] = 2; res.push_back(u); st.pop_back(); }
        }
    }
    return res;
}

int main() {
    // the tree T: 1-2, 1-3, 2-4, 2-5, 3-6, 3-7, 6-8, 6-9 (neighbours in increasing order)
    int n = 9;
    vector<vector<int>> adj(n + 1);
    int E[8][2] = {{1, 2}, {1, 3}, {2, 4}, {2, 5}, {3, 6}, {3, 7}, {6, 8}, {6, 9}};
    for (auto& e : E) { adj[e[0]].push_back(e[1]); adj[e[1]].push_back(e[0]); }
    Timestamps T(adj, 1);
    for (int u = 1; u <= n; u++)
        cout << u << ": tin " << T.tin[u] << " tout " << T.tout[u] << " pos " << T.pos[u]
             << " sz " << T.sz[u] << " dep " << T.dep[u] << "\n";
    cout << "pre-order :"; for (int v : T.order) cout << " " << v; cout << "\n";
    cout << "post-order:"; for (int v : T.post) cout << " " << v; cout << "\n";
    cout << T.isAnc(3, 9) << T.isAnc(2, 6) << T.isAnc(6, 6) << "\n";              // 101
    cout << T.kth(3, 3) << " " << T.kth(3, 5) << " " << T.kth(2, 4) << "\n";        // 8 7 -1
    cout << T.countAtDepth(1, 2) << " " << T.countAtDepth(3, 2) << " "
         << T.countAtDepth(3, 3) << " " << T.countAtDepth(4, 1) << "\n";            // 4 2 2 0
    // a DAG: 1->2, 1->3, 2->4, 3->4, 4->5, 6->3
    vector<vector<int>> dag(7);
    dag[1] = {2, 3}; dag[2] = {4}; dag[3] = {4}; dag[4] = {5}; dag[6] = {3};
    vector<int> f = finishOrder(dag);
    cout << "finish:"; for (int v : f) cout << " " << v; cout << "\n";
    reverse(f.begin(), f.end());
    cout << "topo  :"; for (int v : f) cout << " " << v; cout << "\n";
}
\end{lstlisting}

\begin{lstlisting}[style=trace]
output: 1: tin 0 tout 17 pos 0 sz 9 dep 0      6: tin 8 tout 13 pos 5 sz 3 dep 2
        2: tin 1 tout 6 pos 1 sz 3 dep 1       7: tin 14 tout 15 pos 8 sz 1 dep 2
        3: tin 7 tout 16 pos 4 sz 5 dep 1      8: tin 9 tout 10 pos 6 sz 1 dep 3
        4: tin 2 tout 3 pos 2 sz 1 dep 2       9: tin 11 tout 12 pos 7 sz 1 dep 3
        5: tin 4 tout 5 pos 3 sz 1 dep 2       (one per line in the real output)
        pre-order : 1 2 4 5 3 6 8 9 7
        post-order: 4 5 2 8 9 6 7 3 1
        101
        8 7 -1
        4 2 2 0
        finish: 5 4 2 3 1 6
        topo  : 6 1 3 2 4 5
\end{lstlisting}
On a path of $2\cdot10^5$ vertices the iterative constructor runs without any stack problem; the recursive version of the same code crashes at that depth.

\begin{insightbox}[Which template when?]
\renewcommand{\arraystretch}{1.25}
\begin{tabularx}{\linewidth}{@{}>{\bfseries}l X@{}}
\texttt{isAnc} & ``is $u$ above $v$'', ``does the root path of $v$ pass $u$'', ``which child of $y$ leads to $x$'' \\
\texttt{kth} / \texttt{order} & ``$k$-th vertex of a DFS from $u$'', ``list / iterate the subtree of $u$'', any per-subtree array question (range queries on it: Chapter~15) \\
\texttt{countAtDepth} & ``how many descendants of $u$ in group $g$'' online; offline, the counter difference is shorter \\
\texttt{post} / \texttt{finishOrder} & peel vertices keeping the rest connected; reverse it for a topological order of a DAG \\
\end{tabularx}
\end{insightbox}

%------------------------------------------------------------
\subsubsection*{Common mistakes}

\begin{warningbox}[Mistake 1 --- mixing the two conventions]
\begin{lstlisting}[style=snippet]
int kth = order[tin[u] + k - 1];       // WRONG: tin uses one clock for entry AND exit (0..2n-1)
int kth = order[pos[u] + k - 1];       // RIGHT: pos counts entries only (0..n-1)
\end{lstlisting}
With one clock the subtree is \emph{inside} the interval but does not \emph{fill} it: on $T$, $\tin{3} = 7$ while $\mathrm{pos}[3] = 4$.
\end{warningbox}

\begin{warningbox}[Mistake 2 --- unsorted children when the order matters]
If the problem prescribes ``children by increasing index'' and the adjacency lists were filled from unordered edges, the pre-order --- and every \texttt{kth} answer --- is wrong, while ancestor tests and counts stay right. Sort the lists (or build them sorted).
\end{warningbox}

\begin{warningbox}[Mistake 3 --- binary search on the wrong side of the window]
\begin{lstlisting}[style=snippet]
lower_bound(L, tout[u]) - lower_bound(L, tin[u])      // WRONG only if tout values are in L
upper_bound(L, tout[u]) - lower_bound(L, tin[u])      // RIGHT and robust
\end{lstlisting}
With one clock no entry time equals $\tout{u}$, so both forms happen to agree --- but with the pre-order convention ($[\mathrm{pos}[u], \mathrm{pos}[u]+\mathrm{sz}[u]-1]$) the last position \emph{is} an entry and \texttt{lower\_bound} on the right end drops it. Use the half-open form consistently.
\end{warningbox}

\begin{warningbox}[Mistake 4 --- recursion on a path]
$n = 2\cdot10^5$ with $p_i = i - 1$ is a legal input of all three solved problems; a grid of $500 \times 500$ can force a DFS $125{,}000$ cells deep. Use the iterative template.
\end{warningbox}

\begin{warningbox}[Mistake 5 --- removing in entry order]
The pre-order is the order for \emph{keeping} (every prefix is connected), the post-order the order for \emph{removing} (every prefix is a union of subtrees). Removing the first $k$ vertices of the pre-order removes the root first and can cut the graph into pieces.
\end{warningbox}

\begin{checklistbox}[``I'm stuck'' checklist for Pattern 3]
\begin{itemize}[leftmargin=*]
    \item[$\square$] Can the question be phrased as ``$v$ in the subtree of $u$'' plus something simple? Then the interval $[\tin{u}, \tout{u}]$ or the block $\mathrm{pos}[u]..\mathrm{pos}[u]+\mathrm{sz}[u]-1$ is the answer's skeleton.
    \item[$\square$] Is there a second condition (depth, colour, label)? One sorted list of times per group, or an offline counter difference.
    \item[$\square$] Does the answer depend on the visiting order (``smallest index first'')? Sort the children.
    \item[$\square$] Is it a peeling or dismantling question? Finishing order removes, entry or BFS order keeps.
    \item[$\square$] Where is the root? Re-root at the vertex the question is ``seen from''.
    \item[$\square$] Depth up to $n$? Iterative DFS. Products of sizes? \texttt{long long}.
\end{itemize}
\end{checklistbox}
